%=============================================================
% Paper 3: Campbell & Fisher (2000), AER
%=============================================================
\chapter{Campbell \& Fisher (2000)}

\begin{paperbox}{Aggregate Employment Fluctuations with Microeconomic Asymmetries}
\textbf{Authors:} Jeffrey R.\ Campbell \& Jonas D.M.\ Fisher\\[3pt]
\textbf{Journal:} \textit{American Economic Review} 90(5): 1323--1345\\[3pt]
\textbf{Year:} 2000 \quad\textbf{Field:} Macroeconomics; Labour economics\\[3pt]
\textbf{Classification:} Structural (calibrated dynamic model + simulation)
\end{paperbox}

%------------------------------------------------------------
\section{Research Question}

Can proportional plant-level costs of creating and destroying jobs---net
adjustment costs---explain the empirically documented fact that the job
destruction rate fluctuates more than the job creation rate in U.S.\
manufacturing?

%------------------------------------------------------------
\section{Motivation}

Davis et al.\ (1996) and Davis--Haltiwanger (1992) document that job
destruction is substantially more volatile over the business cycle than job
creation in U.S.\ manufacturing. This asymmetry challenges standard
macroeconomic theory: the one-sector stochastic growth model with homogeneous
plants predicts that job creation and destruction respond equally and
symmetrically to aggregate shocks. Existing explanations either introduce
arbitrary asymmetries in the driving process (Caballero 1992) or are
quantitatively insufficient (Foote 1998). Campbell and Fisher ask whether a
minimal, internally motivated modification---proportional \emph{net} adjustment
costs---can generate the observed asymmetry using a symmetrically distributed
aggregate shock.

%------------------------------------------------------------
\section{Main Contribution}

The paper establishes that proportional net adjustment costs create an
endogenous \textbf{policy asymmetry}: the employment decisions of job-destroying
(shrinking) plants are more sensitive to wage changes than those of
job-creating (expanding) plants. In a calibrated model with continuous
idiosyncratic productivity shocks and ongoing aggregate wage fluctuations,
this microeconomic asymmetry is preserved in the aggregate under plausible
parameter values, generating job destruction volatility that exceeds job
creation volatility---consistent with U.S.\ manufacturing data for the
transportation equipment industry (SIC~37). The contribution is both
analytical (establishing the mechanism) and quantitative (showing the
aggregation condition under which it is preserved).

%------------------------------------------------------------
\section{Relationship to the Literature}

The paper sits at the intersection of the gross job flows literature
(Davis--Haltiwanger 1990, 1992; Mortensen--Pissarides 1994;
Caballero--Hammour 1994, 1996) and the employment adjustment cost literature
(Sargent 1978; Hamermesh 1989; Bentolila--Bertola 1990; Hopenhayn--Rogerson
1993). Its distinctive feature is the emphasis on \emph{net} adjustment costs:
costs that depend on the \emph{number} of jobs created or destroyed, not on
the identity of the workers involved. Campbell and Fisher show that
proportional (rather than quadratic or kinked) adjustment costs generate an
asymmetry even with a symmetric driving process---a result unavailable from
Caballero's (1992) (S,\,s) framework.

%------------------------------------------------------------
\section{Theoretical Framework and Economic Mechanism}

\subsection*{Plant's Problem}

Each plant $i$ produces output $z_t n_t^\alpha$ where $z_t$ is idiosyncratic
productivity (log-random-walk with i.i.d.\ innovation
$\varepsilon_t\sim N(\mu,\sigma_z^2)$) and $n_t$ is employment. The plant
pays net adjustment costs: $t_c$ units of lost output per job added and
$t_d$ units per job subtracted. Defining scaled variables
$x = n_{t-1}/z^{1/(1-\alpha)}$ and $y = n_t/z^{1/(1-\alpha)}$, the Bellman
equation is:
\begin{equation}
  v(x,W) = \max_y\!\left\{
    y^\alpha - Wy - \tau(y,x)(y-x)
    + \beta\,\mathbb{E}\!\left[u'\,v(y/u',W')\mid W\right]
  \right\},
  \label{eq:cf_bellman}
\end{equation}
where $\tau(y,x) = t_c\,\mathbf{1}\{y>x\} - t_d\,\mathbf{1}\{y<x\}$ and
$u' = (z'/z)^{1/(1-\alpha)}$.

\subsection*{Optimal Employment Policy}

The solution defines two scaled thresholds, $y^-(W) < y^+(W)$:
\begin{align}
  \text{Create jobs:} &\quad n^* = y^-(W)\,z^{1/(1-\alpha)}
    &&\text{if } n_{t-1} < n^-(z,W) \\
  \text{Inaction:}    &\quad n^* = n_{t-1}
    &&\text{if } n^-(z,W)\leq n_{t-1}\leq n^+(z,W) \\
  \text{Destroy jobs:}&\quad n^* = y^+(W)\,z^{1/(1-\alpha)}
    &&\text{if } n_{t-1} > n^+(z,W)
\end{align}

\subsection*{Policy Asymmetry}

The elasticities of the scaled employment thresholds with respect to $W$ are:
\begin{align}
  \frac{\partial \ln y^-(W)}{\partial \ln W} &=
    \frac{-1}{1-\alpha}\cdot
    \frac{W}{W + \bigl[1-\beta(1-p)\bigr]t_c + \beta p\,t_d},
  \label{eq:elasticity_create} \\[6pt]
  \frac{\partial \ln y^+(W)}{\partial \ln W} &=
    \frac{-1}{1-\alpha}\cdot
    \frac{W}{W - \bigl[1-\beta(1-p)\bigr]t_d - \beta p\,t_c},
  \label{eq:elasticity_destroy}
\end{align}
where $p$ is the probability of idiosyncratic-shock reversal. Because the
denominator in (\ref{eq:elasticity_destroy}) \emph{subtracts} the adjustment
cost terms while (\ref{eq:elasticity_create})'s denominator \emph{adds} them,
$\left|\partial\ln y^+/\partial\ln W\right| > \left|\partial\ln y^-/\partial\ln
W\right|$ for any strictly positive adjustment costs, regardless of their
relative magnitudes.

\textbf{Intuition:} The total cost of the marginal job for a creating plant is
$W + \text{(creation cost)} + \beta \cdot \text{(expected future destruction cost)}$,
which is \emph{less} elastic to $W$ than $W$ alone. The total cost to a
destroying plant of retaining the marginal job is $W - \text{(avoided
destruction cost)} - \beta\cdot\text{(avoided future creation cost)}$, which
is \emph{more} elastic to $W$. Hence, shrinking plants respond more strongly to
any given wage change.

\subsection*{Aggregation and the Growth-Rate Asymmetry}

The DH gross job flow formulas are:
\begin{align}
  \text{POS}_t &= \frac{2N_{t-1}}{N_t + N_{t-1}}
    \int_0^1 \frac{n_{t-1}(i)}{N_{t-1}}\, g_t(i)^+ \,di, \\[4pt]
  \text{NEG}_t &= \frac{2N_{t-1}}{N_t + N_{t-1}}
    \int_0^1 \frac{n_{t-1}(i)}{N_{t-1}}\, |g_t(i)^-| \,di.
\end{align}
A \textbf{growth-rate asymmetry} works \emph{against} the policy asymmetry:
because $\exp(\cdot)$ is convex, plants near the creation margin respond more
strongly in actual growth rate terms than plants near the destruction margin,
even when policy is symmetric in log terms. The net outcome depends on which
force dominates---and policy asymmetry dominates when idiosyncratic shocks are
\emph{less} persistent.

%------------------------------------------------------------
\section{Data}

\begin{itemize}
  \item \textbf{Calibration target:} Quarterly gross job creation and
    destruction data for SIC~37 (Transportation Equipment) from the LRD/BED.
  \item \textbf{Key moments:} Average net employment growth ($-0.11\%$/quarter),
    average job reallocation rate (11.3\%), standard deviation of employment
    growth (3.55\%), first-order autocorrelation of employment growth (0.08).
  \item \textbf{Calibrated parameters:} $\alpha=0.66$, $\beta=1.052^{1/4}$
    (5\% annual discount); $\sigma_z$ matched to job reallocation rate;
    aggregate shock parameters $(\rho_W, \sigma_W)$ matched to employment
    growth moments.
\end{itemize}

%------------------------------------------------------------
\section{Empirical Strategy}

Calibration and model simulation, not estimation. The strategy:
\begin{enumerate}
  \item Parameterise the model from standard values and calibrated moments
    for SIC~37.
  \item Simulate the model to compute time series of aggregate POS and NEG.
  \item Compare simulated moments (means, standard deviations, covariances of
    POS, NEG, NET) to their empirical counterparts.
  \item Vary adjustment cost parameters $(t_c, t_d)$ within plausible empirical
    bounds to assess sensitivity.
\end{enumerate}

%------------------------------------------------------------
\section{Identification Strategy}

No formal identification in the econometric sense. Predictions are compared to
data descriptively. The identifying variation is the parameterisation of
adjustment costs, taken from external evidence: Los Angeles employer surveys,
Dutch firm studies (Hamermesh et al.\ 1994), Spanish structural estimates
(Aguirregabiria--Alonso-Borrego 1998), and anecdotal evidence from aerospace
production (Benkard 1998). The paper explicitly acknowledges that net
adjustment costs are intrinsically difficult to measure, and that firm surveys
typically measure \emph{gross} adjustment costs (dependent on worker identity)
rather than the \emph{net} adjustment costs (dependent on number of positions)
central to this model.

%------------------------------------------------------------
\section{Main Results}

\subsection*{Analytical Results (Section~I)}

\begin{itemize}
  \item Policy asymmetry holds for \emph{any} strictly positive adjustment
    costs, regardless of the relative magnitudes of $t_c$ and $t_d$.
  \item The asymmetry is stronger when idiosyncratic shocks are \emph{less}
    persistent (smaller $\rho_\text{idio}$), because the policy asymmetry
    dominates the growth-rate asymmetry.
  \item The growth-rate asymmetry always works against the policy asymmetry
    in computing aggregate flows.
\end{itemize}

\subsection*{Quantitative Results (Section~III)}

\begin{center}
\begin{tabular}{lcc}
\toprule
\textbf{Model specification} &
  $\text{std}(\text{NEG})/\text{std}(\text{POS})$ &
  \textbf{Match}? \\
\midrule
No adjustment costs ($t_c=t_d=0$) & $\approx 1.0$ & No \\
Baseline ($t_c=t_d=0.5$, $\rho_W=0.08$) & $\approx 1.8$--$2.0$ & Yes \\
High persistence ($\rho_W=0.5$) & $<1.0$ in some specs & No \\
Empirical SIC~37 & $\approx 1.5$--$2.0$ & --- \\
\bottomrule
\end{tabular}
\end{center}

The model with baseline calibration generates a $\text{std}(\text{NEG})/
\text{std}(\text{POS})$ ratio comparable to the data. The cross-correlation
$\rho(\text{POS},\text{NEG})$ is negative and $\rho(\text{NET},\text{SUM})$
is near zero, broadly consistent with data.

%------------------------------------------------------------
\section{Economic Magnitude and Interpretation}

The SIC~37 industry exhibits $\text{std}(\text{NEG})\approx 2.0$--$2.5\%$,
$\text{std}(\text{POS})\approx 1.0$--$1.5\%$---a ratio of approximately
1.5--2.0. The model with plausible adjustment costs (roughly 0.5--1.5
quarterly wages) replicates this ratio. The mechanism operates through
\emph{aggregate} wage fluctuations, not idiosyncratic shocks---which by
themselves generate symmetric flows. This is an important distinction: the
paper identifies the aggregate component of gross job flow asymmetries.

%------------------------------------------------------------
\section{Mechanisms}

Three forces interact:
\begin{enumerate}
  \item \textbf{Policy asymmetry (main driver):} Adjustment costs make the
    destruction threshold more elastic to wage changes than the creation
    threshold. Each 1\% increase in the wage causes shrinking plants to shrink
    by more than expanding plants expand.
  \item \textbf{Growth-rate asymmetry (countervailing):} Convexity of actual
    vs.\ logarithmic growth rates makes job creators at the margin more
    responsive in actual terms, working against the policy asymmetry.
  \item \textbf{Plant distribution dynamics:} Aggregation over the cross-section
    reflects the endogenous evolution of the distribution of scaled employment
    $x$ across plants. With random-walk idiosyncratic shocks and proportional
    adjustment costs, this evolution preserves (rather than offsets) the policy
    asymmetry.
\end{enumerate}

%------------------------------------------------------------
\section{Robustness Checks}

\begin{itemize}
  \item Results reported for six different combinations of adjustment cost
    parameters and aggregate shock persistence.
  \item Model compared to economy-wide gross job flow data (not just SIC~37)
    with qualitatively similar results.
  \item General equilibrium version of the model developed (Section~II.D)
    to show that the industry-equilibrium results have a valid GE
    interpretation, with the product wage driven by both labour supply and
    demand shocks.
  \item The paper explicitly notes that \emph{gross} adjustment costs (costs
    that depend on worker identity, not net employment changes) would
    \emph{not} generate the policy asymmetry, providing a discriminating
    prediction of the net vs.\ gross cost distinction.
\end{itemize}

%------------------------------------------------------------
\section{Limitations}

\begin{enumerate}
  \item \textbf{No matching frictions:} Abstracts entirely from search and
    matching; the Mortensen--Pissarides mechanism (asymmetric matching vs.\
    separation thresholds) could compound or offset the policy asymmetry.
  \item \textbf{Symmetric aggregate shocks:} Assumes symmetrically distributed
    aggregate uncertainty; in reality, recessions may be more severe than
    expansions, compounding the asymmetry.
  \item \textbf{Calibration, not estimation:} Adjustment costs are selected
    from a plausible range, not formally estimated. The model cannot be
    rejected on these dimensions.
  \item \textbf{Manufacturing only:} Service sectors, where labour market
    dynamics may differ substantially, are excluded.
  \item \textbf{No entry/exit:} The population of plants is fixed, abstracting
    from births and deaths that account for 20--25\% of gross job flows in
    DH1992.
\end{enumerate}

%------------------------------------------------------------
\section{Policy Implications}

\begin{itemize}
  \item Firing costs (higher $t_d$) \emph{reduce} employment volatility by
    dampening the destruction response, consistent with the Bentolila--Bertola
    (1990) motivation for European firing-cost policies---but at the cost of
    reduced flexibility.
  \item Policies reducing hiring costs specifically for expanding plants would
    raise the creation elasticity and reduce the volatility asymmetry.
  \item The paper provides a new rationale for the excess volatility of
    unemployment in recessions: it partly reflects the greater sensitivity of
    plant-level employment cuts to aggregate shocks, not just ``cleansing''
    of less productive matches.
\end{itemize}

%------------------------------------------------------------
\section{Overall Assessment}

Campbell and Fisher (2000) is a rigorous paper that makes a genuinely original
theoretical point---the policy asymmetry mechanism---and shows it has
quantitative bite. The general equilibrium embedding is careful. The main
weakness is the calibration approach, which limits the paper's ability to
formally assess how much of the observed asymmetry is explained by the
mechanism. The paper is less influential than DH1992 partly because the
Mortensen--Pissarides search-and-matching framework became dominant in the
subsequent decade, and partly because the mechanism requires aggregate shocks
to operate primarily through wages.

\textbf{Quality: Good. A solid theoretical contribution with careful
quantitative assessment.}

%------------------------------------------------------------
\section{Relevance to My Research}

The policy asymmetry mechanism has direct relevance for understanding how
Canadian firms respond to aggregate shocks (commodity price cycles, exchange
rate changes) in different ways depending on whether they are expanding or
contracting. The key testable prediction: the elasticity of employment cuts to
adverse shocks should exceed the elasticity of employment additions to
favourable shocks of the same magnitude, and this asymmetry should be larger
for less persistent aggregate shocks (e.g., commodity price swings).

%------------------------------------------------------------
\section{Potential Research Extensions}

\begin{itemize}
  \item Estimate $t_c$ and $t_d$ structurally from Canadian establishment-level
    data by matching moments of the job creation/destruction distribution,
    allowing estimation of the model's degree of asymmetry.
  \item Test the model's prediction that industries with more transient aggregate
    shocks (e.g., commodity sectors) show larger $\text{std}(\text{NEG})/
    \text{std}(\text{POS})$ ratios than industries with persistent shocks.
\end{itemize}

%------------------------------------------------------------
\section{Research Gaps}

\begin{itemize}
  \item The paper does not model worker heterogeneity or worker-level welfare
    consequences of the asymmetric adjustment process.
  \item There is no explanation for why $t_c$ and $t_d$ should differ, or
    how they vary across industries, countries, or regulatory regimes.
\end{itemize}

%------------------------------------------------------------
\section{Methodological Lessons}

\begin{enumerate}
  \item The DH growth rate metric versus the log-difference matters for
    quantitative model evaluation: the growth-rate asymmetry is a real force
    that works against the policy asymmetry and would be missed using
    log-approximations.
  \item The policy asymmetry result holds \emph{analytically} for any positive
    adjustment costs---showing a qualitative result before turning to
    quantitative assessment is methodologically important.
  \item Comparing model-generated statistics to industry-level data (SIC~37)
    rather than economy-wide data is the right level of aggregation when the
    model is an industry-equilibrium model.
\end{enumerate}

%------------------------------------------------------------
\section*{Five-Sentence Summary}

Campbell and Fisher (2000) show that proportional net adjustment costs create
a policy asymmetry at the plant level: shrinking plants are more sensitive to
wage changes than expanding plants, because adjustment costs add to the total
cost of job creation but reduce the total cost of job retention. This
microeconomic asymmetry, when preserved through aggregation, generates a model
in which the rate of job destruction fluctuates more than the rate of job
creation in response to aggregate shocks---a key stylised fact from U.S.\
manufacturing that standard representative-agent models fail to reproduce. A
growth-rate asymmetry (convexity of actual vs.\ log growth rates) works in the
opposite direction and can dominate when idiosyncratic shocks are highly
persistent, so the mechanism is strongest for relatively transient aggregate
fluctuations. A calibrated version of the model for the U.S.\ transportation
equipment industry generates a ratio $\text{std}(\text{NEG})/\text{std}(
\text{POS})$ comparable to the data, explaining most of the observed excess
volatility of job destruction without requiring an asymmetrically distributed
driving process. The paper provides a clean, internally motivated explanation
for job flow asymmetries based on the loss of production incurred when
reorganising a plant to operate at a larger or smaller scale.

\vspace{6pt}
\begin{quickref}
\begin{tabularx}{\linewidth}{@{}lX@{}}
\textbf{Key contribution} &
  Analytical policy asymmetry mechanism: net adjustment costs make job
  destruction more elastic to aggregate shocks than job creation, even with
  a symmetric driving process. \\[4pt]
\textbf{Key weakness} &
  Calibration (not estimation) of adjustment costs; no entry/exit; model
  cannot be formally tested against alternative explanations. \\[4pt]
\textbf{Most interesting gap} &
  Does the policy asymmetry vary systematically across industries with
  different regulatory environments (high vs.\ low firing costs)? A
  cross-country test of the mechanism would be revealing. \\[4pt]
\textbf{Publishable extension} &
  Test the Campbell--Fisher mechanism using matched employer-employee data
  from Canada, comparing employment creation and destruction elasticities
  to province-level commodity price cycles across industries with high and
  low regulatory adjustment cost environments.
\end{tabularx}
\end{quickref}
